\section{Introduction et Étude Préalable}

%------------------------------------------------------------
% Slide 3 : Contexte du Projet
%------------------------------------------------------------
\begin{frame}{Contexte du Projet : DataLab Research}
    \begin{columns}[T]
        \begin{column}{0.5\textwidth}
            \textbf{L'Organisation :}
            \begin{itemize}
                \item Institut de recherche scientifique moderne.
                \item \textbf{80 Utilisateurs permanents} (Chercheurs, Ingénieurs, Administrateurs).
                \item Manipulation quotidienne de volumes massifs de données scientifiques sensibles.
            \end{itemize}
        \end{column}
        
        \begin{column}{0.5\textwidth}
            \textbf{Exigences Fondamentales :}
            \begin{block}{Haute Disponibilité (24/7)}
                Les projets de recherche en cours ne peuvent souffrir d'aucune interruption de service.
            \end{block}
            \begin{block}{Sécurité Absolue}
                Garantir la confidentialité et l'étanchéité des flux face aux accès non autorisés.
            \end{block}
        \end{column}
    \end{columns}
\end{frame}

%------------------------------------------------------------
% Slide 4 : Étude de l'Existant et Problématique
%------------------------------------------------------------
\begin{frame}{Étude de l'Existant et Problématique}
    \begin{alertblock}{Limitation de l'Infrastructure d'Origine}
        Absence de segmentation réseau, absence de redondance sur la couche de stockage et maintenance purement réactive (faute d'outils de monitoring).
    \end{alertblock}

    \vspace{0.3cm}
    \textbf{La Problématique Centrale :}
    \begin{exampleblock}{\textit{Question d'Ingénierie}}
        Comment bâtir une infrastructure datacenter virtualisée, hautement disponible et sécurisée, capable d'isoler les flux applicatifs tout en offrant une observabilité totale et automatisée aux administrateurs ?
    \end{exampleblock}
\end{frame}

%------------------------------------------------------------
% Slide 5 : Démarche Méthodologique et Contraintes
%------------------------------------------------------------
\begin{frame}{Démarche Méthodologique et Contraintes}
    Face à des contraintes fortes, notre équipe a adopté une démarche d'ingénierie pragmatique.

    \begin{columns}[T]
        \begin{column}{0.48\textwidth}
            \begin{block}{Contraintes Réelles}
                \begin{itemize}
                    \item \textbf{Temps Restreint :} Calendrier global non en notre faveur.
                    \item \textbf{Matériel Limité :} Postes de simulation restreints en CPU/RAM.
                \end{itemize}
            \end{block}
        \end{column}

        \begin{column}{0.48\textwidth}
            \begin{block}{La Solution : Approche MVP}
                \begin{itemize}
                    \item \textbf{Conception complète :} À l'état de l'art (8 VLANs théoriques).
                    \item \textbf{Déploiement ciblé :} Mise en œuvre d'un MVP fonctionnel  pour démontrer la faisabilité et la valeur ajoutée.
                \end{itemize}
            \end{block}
        \end{column}
    \end{columns}
\end{frame}

%------------------------------------------------------------
% Slide 6 : Gestion de Projet et Outils Collaboratifs
%------------------------------------------------------------
\begin{frame}{Gestion de Projet et Outils Collaboratifs}
    Pour orchestrer le travail d'équipe sous un cycle court, deux outils piliers ont été adoptés :

    \begin{columns}[T]
        \begin{column}{0.5\textwidth}
            \begin{itemize}
                \item \textbf{Suivi  (Notion) :}
                Centralisation de la phase de conception, répartition des rôles et suivi visuel des tâches.
                \item \textbf{Versionnage (GitHub) :}
                Source unique de vérité pour l'ensemble des scripts de configuration réseau et système.
            \end{itemize}
            \vspace{0.2cm}
            \centering
            \tiny{\url{https://github.com/Jovialngandu/datalab_infra_m1}}
        \end{column}

        \begin{column}{0.5\textwidth}
            \begin{figure}
                \centering
                % Intégration de votre capture réelle notion.png
                \includegraphics[width=0.95\textwidth]{figures/notion.png}
                \caption{Tableau  de l'équipe sur Notion}
            \end{figure}
        \end{column}
    \end{columns}
\end{frame}